\documentclass[12pt, a4paper]{article}

\usepackage[margin=1in]{geometry}
\usepackage{amsmath}
\usepackage{amssymb}
\usepackage{graphicx}
\usepackage{booktabs}
\usepackage{array}
\usepackage{multirow}
\usepackage{hyperref}
\usepackage{xcolor}
\usepackage{enumitem}
\usepackage{caption}
\usepackage{subcaption}
\usepackage{setspace}
\usepackage{parskip}
\usepackage{titlesec}
\usepackage{fancyhdr}
\usepackage{cite}

\hypersetup{
    colorlinks=true,
    linkcolor=blue!70!black,
    citecolor=green!50!black,
    urlcolor=blue!70!black,
}

\titleformat{\section}{\large\bfseries}{\thesection.}{0.5em}{}
\titleformat{\subsection}{\normalsize\bfseries}{\thesubsection}{0.5em}{}

\pagestyle{fancy}
\fancyhf{}
\rhead{\small CS220: Computer Vision Project Report}
\lhead{\small Image Retrieval Using Color and Texture Features}
\rfoot{\thepage}
\renewcommand{\headrulewidth}{0.4pt}

\setlength{\parskip}{6pt}
\onehalfspacing

\begin{document}

% ─── Title ───────────────────────────────────────────────────────────────────
\begin{titlepage}
    \centering
    \vspace*{2cm}
    {\LARGE \textbf{Image Retrieval Using Color and Texture Features}}\\[1em]
    {\large A Content-Based Image Retrieval System}\\[2.5em]

    {\normalsize
        Atharva Nanoti (220001013) \quad
        Shaurya Khetarpal (220001074) \quad
        Saket Meshram (220001066)
    }\\[0.5em]
    {\small Indian Institute of Technology Indore}\\[2em]
    {\small \today}\\[2em]

    \begin{abstract}
        \noindent
        This report presents the design and evaluation of a Content-Based Image Retrieval
        (CBIR) system that retrieves visually similar images purely from their pixel-level
        content. Three categories of hand-crafted visual descriptors are implemented: a
        3-channel color histogram, an optimized Local Binary Pattern (LBP) descriptor, and
        statistical features derived from the Gray-Level Co-occurrence Matrix (GLCM).
        All combinations of these descriptors are also evaluated as fused feature vectors.
        Retrieval performance is measured via Precision@$K$ ($K \in \{1,5,10,20\}$) on the
        Wang dataset (1\,000 images, 10 classes), using every image as a query in a
        leave-one-out protocol. Results show that fusing all three descriptors achieves
        the best global Precision@10 of \textbf{0.6699}, a 7.9\% relative improvement
        over the color histogram alone, while also outperforming every single-descriptor
        baseline.
    \end{abstract}
\end{titlepage}

% ─── 1. Introduction ─────────────────────────────────────────────────────────
\section{Introduction}

The exponential growth of digital image collections—spanning medical imaging, satellite
photography, social media archives, and e-commerce catalogues—has made efficient and
semantically meaningful image search a central problem in computer vision. Traditional
keyword-based retrieval relies on manual annotation, which does not scale and is inherently
subjective. Content-Based Image Retrieval (CBIR) addresses this limitation by describing and
indexing images directly from their visual content, enabling similarity search without any
human-supplied labels \cite{smeulders2000}.

In a CBIR pipeline, each image is mapped to a compact numerical feature vector. A query
image is then processed by the same feature extractor, and database images are ranked by
the distance between their feature vectors and the query vector. The quality of retrieval is
therefore determined primarily by the expressiveness and discriminability of the chosen
feature representation.

Classical hand-crafted descriptors based on color and texture remain widely used in retrieval
benchmarks because they are computationally lightweight, interpretable, and do not require
labelled training data \cite{rui1999}. Color histograms capture the global hue composition of an
image \cite{swain1991}. Local Binary Patterns (LBP) encode micro-texture by recording
binarized pixel-neighborhood relationships \cite{ojala2002}. Gray-Level Co-occurrence
Matrix (GLCM) features capture higher-order statistical co-occurrence structure that is
orthogonal to both color and micro-texture \cite{haralick1973}.

This project builds and evaluates a complete CBIR system using these three descriptor
families, individually and in all pairwise and triple combinations, on the standard Wang
benchmark dataset. The key contributions are:

\begin{itemize}[noitemsep]
    \item A modular, fully cached retrieval pipeline supporting seven distinct feature extractors.
    \item A comprehensive leave-one-out evaluation over all 1\,000 images, reporting
          Precision@$K$ globally and per semantic class.
    \item A systematic comparison revealing the complementarity of color and texture
          information for image retrieval.
\end{itemize}

% ─── 2. Problem Statement ────────────────────────────────────────────────────
\section{Problem Statement}

Given a database $\mathcal{D} = \{I_1, I_2, \ldots, I_N\}$ of $N = 1\,000$ images and a
query image $I_q \in \mathcal{D}$, the goal is to return the $K$ most visually similar images
from $\mathcal{D} \setminus \{I_q\}$ such that images from the same semantic category as
$I_q$ rank as high as possible.

Formally, each image $I_i$ is mapped to a feature vector $\mathbf{f}_i \in \mathbb{R}^d$
by a feature extractor $\phi: \mathcal{I} \rightarrow \mathbb{R}^d$. Retrieval is performed
by ranking all database images by a distance function $\delta(\mathbf{f}_q, \mathbf{f}_i)$ in
ascending order, and returning the top-$K$ results $\mathcal{R}_K(I_q)$.

Performance is quantified by Precision@$K$:
\begin{equation}
    \text{P@}K(I_q) = \frac{1}{K} \sum_{I_i \in \mathcal{R}_K(I_q)}
                      \mathbf{1}\bigl[\text{class}(I_i) = \text{class}(I_q)\bigr]
\end{equation}
where $\mathbf{1}[\cdot]$ is the indicator function and class labels are derived from
filenames (images $\{100i,\ldots,100(i+1)-1\}.jpg$ belong to class $i$).

The core challenge is to identify which feature representation $\phi$ maximises retrieval
precision across diverse visual categories using only unsupervised, pixel-level descriptors
and a fixed Euclidean distance metric.

% ─── 3. Approach / Method ───────────────────────────────────────────────────
\section{Approach and Method}

\subsection{System Pipeline}

The retrieval system follows a six-stage pipeline:

\begin{enumerate}[noitemsep]
    \item \textbf{Preprocessing.} Images are read in BGR format via OpenCV and converted
          to RGB (for color features) or grayscale (for texture features) as required.
    \item \textbf{Feature extraction.} Each image is mapped to a fixed-length feature vector
          by one of the seven extractors described in Section~\ref{sec:extractors}.
    \item \textbf{Database construction.} Feature vectors for all $N$ images are computed
          once and persisted as Python pickle files to avoid redundant computation on
          subsequent runs.
    \item \textbf{Query processing.} A query image's feature vector is computed on the fly
          using the same extractor.
    \item \textbf{Similarity ranking.} Euclidean distance between the query vector and every
          database vector is computed, and images are ranked in ascending order of distance.
    \item \textbf{Evaluation.} The top-$K$ returned images (excluding the query itself) are
          compared against the ground-truth class label of the query to compute Precision@$K$.
\end{enumerate}

\subsection{Feature Extractors}
\label{sec:extractors}

\subsubsection{Color Histogram (CH)}

The color histogram represents the global color distribution of an image. The image is
converted to RGB and a joint 3-channel histogram is computed with $8 \times 8 \times 8$
bins spanning $[0, 255]$ in each channel. The resulting histogram is $\ell^2$-normalized and
flattened to produce a 512-dimensional feature vector.

\begin{equation}
    \mathbf{f}_{\text{CH}} = \text{normalize}\Bigl(
        \text{hist}_{8,8,8}(I_R, I_G, I_B)
    \Bigr) \in \mathbb{R}^{512}
\end{equation}

Color histograms are invariant to spatial rearrangements of pixels and robust to small
geometric distortions, making them effective for categories with distinctive palettes (e.g.
flowers, dinosaurs) but weaker for categories sharing similar color ranges (e.g. beaches and
mountains).

\subsubsection{Local Binary Patterns -- Fast (LBP)}

Local Binary Patterns encode the local texture of an image by comparing each pixel's
intensity with its circularly sampled neighbors \cite{ojala2002}. We use the \emph{uniform}
LBP variant from \texttt{scikit-image}, with $P = 24$ sampling points and radius $R = 3$.
Uniform patterns (at most two 0--1 transitions) are assigned individual bins; all
non-uniform patterns are grouped into a single bin, yielding $P + 2 = 26$ bins. The
resulting histogram is $\ell^1$-normalized.

\begin{equation}
    \text{LBP}_{P,R}(x, y) =
    \sum_{p=0}^{P-1} s\bigl(g_p - g_c\bigr)\,2^p, \quad
    s(u) = \begin{cases} 1 & u \geq 0 \\ 0 & u < 0 \end{cases}
\end{equation}

where $g_c$ is the gray value of the center pixel and $g_p$ are the $P$ equally-spaced
neighbor values on a circle of radius $R$. This yields a 26-dimensional texture descriptor.

\subsubsection{Gray-Level Co-occurrence Matrix Features (GLCM)}

The GLCM captures second-order texture statistics by recording co-occurrence frequencies
of pixel-intensity pairs at a given spatial offset \cite{haralick1973}. We compute the GLCM
at four angular directions ($0°, 45°, 90°, 135°$) with a single pixel distance, using
256 gray levels and symmetric, normalized co-occurrence matrices. Five Haralick statistics
are extracted from each directional matrix:

\begin{itemize}[noitemsep]
    \item \textbf{Contrast}: $\sum_{i,j} (i-j)^2\, p_{ij}$
    \item \textbf{Dissimilarity}: $\sum_{i,j} |i-j|\, p_{ij}$
    \item \textbf{Homogeneity}: $\sum_{i,j} p_{ij} / (1 + (i-j)^2)$
    \item \textbf{Energy}: $\sqrt{\sum_{i,j} p_{ij}^2}$
    \item \textbf{Correlation}: $\sum_{i,j} \dfrac{(i-\mu_i)(j-\mu_j)\,p_{ij}}{\sigma_i\sigma_j}$
\end{itemize}

With 5 statistics across 4 angles, the GLCM descriptor is a 20-dimensional vector.

\subsubsection{Combined Feature Representations}

To exploit the complementarity of color and texture, all combinations of the three base
descriptors are also evaluated. Before concatenation each individual descriptor is
$\ell^2$-normalized to unit length so that no descriptor dominates due to scale differences.
The four combined extractors and their dimensionalities are:

\begin{table}[h]
    \centering
    \caption{Combined feature extractors and their dimensionalities.}
    \label{tab:dims}
    \begin{tabular}{@{}lcc@{}}
        \toprule
        \textbf{Extractor} & \textbf{Components} & \textbf{Dimension} \\
        \midrule
        CH + LBP            & 512 + 26          & 538  \\
        CH + GLCM           & 512 + 20          & 532  \\
        LBP + GLCM          & 26 + 20           & 46   \\
        CH + LBP + GLCM     & 512 + 26 + 20     & 558  \\
        \bottomrule
    \end{tabular}
\end{table}

\subsection{Similarity Measure}

Retrieval is performed using the \textbf{Euclidean distance} between feature vectors:

\begin{equation}
    \delta(\mathbf{f}_q, \mathbf{f}_i) = \|\mathbf{f}_q - \mathbf{f}_i\|_2
    = \sqrt{\sum_{k=1}^{d} (f_{q,k} - f_{i,k})^2}
\end{equation}

All database images are ranked by this distance in ascending order, and the query image
itself (which achieves distance zero when it appears in the database) is excluded from the
returned results before precision is computed.

\subsection{Evaluation Protocol}

Performance is measured using a \textbf{leave-one-out} protocol: each of the 1\,000 images
in the Wang dataset is used as a query in turn. For each query, the top-$K$ retrieved images
(with the query excluded) are compared against the query's ground-truth class label.

Precision@$K$ is computed for $K \in \{1, 5, 10, 20\}$ and reported in two ways:

\begin{itemize}[noitemsep]
    \item \textbf{Global}: macro-average over all 1\,000 queries.
    \item \textbf{Per-class}: macro-average over the 100 queries belonging to each class.
\end{itemize}

% ─── 4. Results and Analysis ─────────────────────────────────────────────────
\section{Results and Analysis}

\subsection{Global Precision@K}

Table~\ref{tab:global} reports the global Precision@$K$ for all seven feature extractors,
averaged over all 1\,000 queries. Several clear trends emerge.

\begin{table}[h]
    \centering
    \caption{Global Precision@$K$ averaged over all 1\,000 queries on the Wang dataset.
             Best values in each column are \textbf{bold}.}
    \label{tab:global}
    \setlength{\tabcolsep}{8pt}
    \begin{tabular}{@{}lcccc@{}}
        \toprule
        \textbf{Feature Extractor} & \textbf{P@1} & \textbf{P@5} & \textbf{P@10} & \textbf{P@20} \\
        \midrule
        \multicolumn{5}{@{}l}{\textit{Individual descriptors}} \\
        Color Histogram (CH)      & 0.7180 & 0.6668 & 0.6210 & 0.5693 \\
        LBP Fast                  & 0.5650 & 0.5150 & 0.4893 & 0.4409 \\
        GLCM Features             & 0.4130 & 0.3594 & 0.3339 & 0.3080 \\
        \midrule
        \multicolumn{5}{@{}l}{\textit{Combined descriptors}} \\
        CH + LBP                  & 0.7610 & 0.7058 & 0.6629 & 0.6134 \\
        CH + GLCM                 & 0.7290 & 0.6730 & 0.6287 & 0.5781 \\
        LBP + GLCM                & 0.6040 & 0.5390 & 0.5044 & 0.4677 \\
        \textbf{CH + LBP + GLCM} & \textbf{0.7660} & \textbf{0.7120} & \textbf{0.6699} & \textbf{0.6172} \\
        \bottomrule
    \end{tabular}
\end{table}

\textbf{Color dominates texture.} Among individual descriptors, the color histogram achieves
the highest retrieval precision across all values of $K$ (P@10 = 0.6210), substantially
outperforming both LBP (P@10 = 0.4893) and GLCM (P@10 = 0.3339). This is consistent
with the nature of the Wang dataset, where many categories (flowers, dinosaurs, horses,
beaches) are visually distinguishable by their dominant color distributions.

\textbf{GLCM is the weakest individual descriptor.} Despite capturing second-order spatial
texture structure, GLCM features alone yield the lowest precision. This is partly because a
20-dimensional vector carries limited discriminative capacity for a 10-class dataset, and
because structural texture alone is insufficient to distinguish categories with similar gray-level
co-occurrence patterns (e.g. buildings and mountains).

\textbf{Feature fusion consistently improves retrieval.} Every combined extractor outperforms
its constituent individual descriptors. Combining CH with LBP yields a relative improvement of
+6.7\% over CH alone at P@10 (0.6629 vs 0.6210). Adding GLCM on top further raises
precision to 0.6699. This demonstrates that color and texture information are complementary:
texture features help distinguish categories that share similar color distributions.

\textbf{The full combination CH + LBP + GLCM is the best extractor overall}, achieving the
highest precision at every $K$ value. The gain is consistent but modest beyond K=10,
suggesting that at larger retrieval depths the additional discriminability from texture
descriptors is partially offset by harder-to-rank within-class variation.

\subsection{Per-Class Precision@10}

Table~\ref{tab:perclass} reports Precision@10 broken down by semantic class for all seven
extractors. This reveals important class-specific patterns.

\begin{table}[h]
    \centering
    \caption{Per-class Precision@10 for all extractors. Classes sorted by difficulty.}
    \label{tab:perclass}
    \setlength{\tabcolsep}{5pt}
    \begin{tabular}{@{}lccccccc@{}}
        \toprule
        \textbf{Class} &
        \textbf{CH} &
        \textbf{LBP} &
        \textbf{GLCM} &
        \textbf{CH+LBP} &
        \textbf{CH+GLCM} &
        \textbf{LBP+GLCM} &
        \textbf{CH+LBP+GLCM} \\
        \midrule
        Dinosaurs    & 0.987 & 0.921 & 0.363 & 0.987 & 0.987 & 0.919 & 0.987 \\
        Horses       & 0.925 & 0.642 & 0.338 & 0.941 & 0.928 & 0.643 & 0.944 \\
        Flowers      & 0.650 & 0.745 & 0.805 & 0.832 & 0.689 & 0.709 & 0.834 \\
        Africans     & 0.713 & 0.412 & 0.183 & 0.736 & 0.718 & 0.386 & 0.738 \\
        Food         & 0.604 & 0.360 & 0.235 & 0.635 & 0.606 & 0.352 & 0.650 \\
        Buses        & 0.533 & 0.785 & 0.592 & 0.654 & 0.557 & 0.760 & 0.672 \\
        Elephants    & 0.610 & 0.291 & 0.226 & 0.626 & 0.615 & 0.370 & 0.636 \\
        Buildings    & 0.465 & 0.267 & 0.231 & 0.486 & 0.460 & 0.349 & 0.498 \\
        Mountains    & 0.359 & 0.225 & 0.160 & 0.367 & 0.355 & 0.252 & 0.361 \\
        Beaches      & 0.364 & 0.245 & 0.206 & 0.365 & 0.372 & 0.304 & 0.379 \\
        \bottomrule
    \end{tabular}
\end{table}

\textbf{Dinosaurs are the easiest class.} All color-based methods achieve P@10 $\approx$
0.987, near-perfect retrieval. Dinosaur images share a highly distinctive combination of
green/grey coloration against uniform backgrounds, making them separable by color histogram
alone. Notably, GLCM performs poorly here (0.363), indicating that texture alone is
insufficient when images have globally similar texture statistics across different classes.

\textbf{Horses are the second easiest.} Color-combination methods exceed 0.94 at P@10.
The brown/chestnut tones of horse bodies combined with grass-green backgrounds create a
strong color signature.

\textbf{Flowers are the one class where GLCM excels.} GLCM achieves P@10 = 0.805 for
flowers, the highest GLCM score across all classes and competitive with LBP (0.745). This
is because flower petals produce distinctive high-contrast, regular co-occurrence patterns
that GLCM captures effectively. It is also the only class where GLCM outperforms the color
histogram (CH P@10 = 0.650).

\textbf{Buses benefit most from LBP.} LBP achieves 0.785 for buses, the highest single-
descriptor score for this class and surpassing the color histogram (0.533). The strong
rectangular structural patterns of bus bodies and windows produce discriminative LBP histograms.

\textbf{Mountains and Beaches are the hardest classes} across all extractors. Mountains
are frequently confused with beaches due to overlapping blue-sky and grey-brown color
distributions. Neither texture descriptor helps significantly: even the best extractor
achieves only P@10 = 0.379 for beaches and 0.367 for mountains. This limitation is
inherent to global feature representations that lack spatial context.

\subsection{Effect of $K$ on Precision}

Across all extractors and classes, Precision@$K$ decreases monotonically as $K$ increases.
This is expected: as more results are retrieved, the fraction of relevant images in the
returned set naturally falls because each class contains only 99 same-class images to
retrieve. The relative ordering of extractors is stable across all $K$ values, confirming
that the conclusions drawn at P@10 generalise to other retrieval depths.

% ─── 5. Future Scope ────────────────────────────────────────────────────────
\section{Future Scope of the Work}

The current system establishes a strong classical baseline but has several avenues for
improvement:

\textbf{Deep feature representations.} Convolutional neural network features (e.g. ResNet,
VGG activations) have consistently outperformed hand-crafted descriptors on image retrieval
benchmarks \cite{babenko2014}. Replacing or augmenting the hand-crafted vectors with
deep features extracted from a pre-trained backbone would likely yield large improvements
particularly for visually ambiguous classes such as beaches and mountains.

\textbf{Alternative distance metrics.} Only Euclidean distance is used in the current
implementation. Exploring Chi-squared distance (well-suited to histograms), cosine similarity,
or learned metric embeddings could improve retrieval precision without changing the feature
representations.

\textbf{Spatial pyramid features.} The current descriptors are computed globally over the
entire image, discarding spatial structure. Spatial pyramid matching \cite{lazebnik2006}
divides images into progressively finer grids and pools descriptors at each level, encoding
spatial layout while retaining some degree of invariance.

\textbf{Relevance feedback.} Interactive retrieval systems allow users to mark retrieved
images as relevant or irrelevant. This feedback can be used to re-weight or shift the query
vector towards the relevant image region in feature space, progressively improving results
without retraining.

\textbf{Approximate nearest neighbour search.} The current retrieval scans all $N$ feature
vectors exhaustively, giving $O(Nd)$ query time. For large-scale databases, locality-sensitive
hashing (LSH) or product quantization can reduce query time to near $O(\log N)$ with minimal
precision loss.

\textbf{Recall and mean Average Precision.} The current evaluation uses only Precision@$K$.
Reporting mean Average Precision (mAP) and recall curves would give a more complete picture
of the precision-recall trade-off across all retrieval depths.

% ─── 6. References ──────────────────────────────────────────────────────────
\section*{References}
\addcontentsline{toc}{section}{References}

\begin{thebibliography}{9}

\bibitem{smeulders2000}
A.~W.~M.~Smeulders, M.~Worring, S.~Santini, A.~Gupta, and R.~Jain,
``Content-based image retrieval at the end of the early years,''
\textit{IEEE Transactions on Pattern Analysis and Machine Intelligence},
vol.~22, no.~12, pp.~1349--1380, Dec.~2000.

\bibitem{rui1999}
Y.~Rui, T.~S.~Huang, and S.-F.~Chang,
``Image retrieval: Current techniques, promising directions, and open issues,''
\textit{Journal of Visual Communication and Image Representation},
vol.~10, no.~1, pp.~39--62, Mar.~1999.

\bibitem{swain1991}
M.~J.~Swain and D.~H.~Ballard,
``Color indexing,''
\textit{International Journal of Computer Vision},
vol.~7, no.~1, pp.~11--32, Nov.~1991.

\bibitem{ojala2002}
T.~Ojala, M.~Pietik\"{a}inen, and T.~M\"{a}enp\"{a}\"{a},
``Multiresolution gray-scale and rotation invariant texture classification
with local binary patterns,''
\textit{IEEE Transactions on Pattern Analysis and Machine Intelligence},
vol.~24, no.~7, pp.~971--987, Jul.~2002.

\bibitem{haralick1973}
R.~M.~Haralick, K.~Shanmugam, and I.~Dinstein,
``Textural features for image classification,''
\textit{IEEE Transactions on Systems, Man, and Cybernetics},
vol.~3, no.~6, pp.~610--621, Nov.~1973.

\bibitem{wang2001}
J.~Li and J.~Z.~Wang,
``Automatic linguistic indexing of pictures by a statistical modeling approach,''
\textit{IEEE Transactions on Pattern Analysis and Machine Intelligence},
vol.~25, no.~9, pp.~1075--1088, Sep.~2003.

\bibitem{babenko2014}
A.~Babenko, A.~Slesarev, A.~Chigorin, and V.~Lempitsky,
``Neural codes for image retrieval,''
in \textit{Proc.\ European Conference on Computer Vision (ECCV)},
pp.~584--599, 2014.

\bibitem{lazebnik2006}
S.~Lazebnik, C.~Schmid, and J.~Ponce,
``Beyond bags of features: Spatial pyramid matching for recognizing natural scene categories,''
in \textit{Proc.\ IEEE Conference on Computer Vision and Pattern Recognition (CVPR)},
pp.~2169--2178, 2006.

\end{thebibliography}

\end{document}
